\begin{table}[!htbp]
\caption{变量、分析单位及预测时可用性。“是，滞后”表示一步超前预测仅使用上一
聚合期的信息。}
\label{tab:variable_availability}
\centering
\footnotesize
\begin{tabularx}{\textwidth}{p{2.3cm}Xp{1.9cm}p{1.9cm}p{2.2cm}}
\toprule
\textbf{变量} & \textbf{含义与作用} & \textbf{空间单位} &
\textbf{时间单位} & \textbf{预测时可用}\\
\midrule
犯罪计数 & 三类统一犯罪的每日事件数；未分箱预测目标 & 1 km 网格 & 日 &
目标，否\\
滞后犯罪 & 上一时期犯罪计数，离散为 \(\{0,1,2,\geq3\}\)；基准状态 &
1 km 网格 & 前一日 & 是，滞后\\
自行车流入 & 已解析公共端点的到达数；移动原始量 & 1 km 网格 & 前一日 &
是，滞后\\
自行车流出 & 已解析公共端点的出发数；移动原始量 & 1 km 网格 & 前一日 &
是，滞后\\
自行车总流量 & 流入与流出之和，按训练期阈值离散；新增数据状态 &
1 km 网格 & 前一日 & 是，滞后\\
日历状态 & 星期、气象季节和公共假日状态；基准状态 & 城市/网格 & 预测日 &
是\\
网格标识 & 验证与测试前固定的空间状态；基准状态 & 1 km 网格 & 固定 & 是\\
\bottomrule
\end{tabularx}
\end{table}
